\section{Topology Modulation Operator}

\subsection{Overview}

The topology modulation operator completes the geometric transformation
pipeline at altitude A8 by incorporating global structural information
from the manifold. Whereas spectral, curvature, and holonomy operators
capture local and accumulated geometric behavior, the topology operator
extracts a global signature that reflects the manifold’s large-scale
connectivity and shape.

This section formalizes the topology signature, its construction, and
its role in the Redux embedding pipeline.

\subsection{Topological Structure of the Manifold}

Although the manifold $M = (\mathbb{R}^{80}, g)$ is topologically
trivial, the metric tensor $g$ induces nontrivial geometric features
that can be summarized through a topology signature. The purpose of the
topology operator is not to detect classical invariants such as Betti
numbers, but rather to extract a global modulation factor derived from
the manifold’s geometric configuration.

\subsection{Topology Signature}

We define the topology signature
\[
\tau(x) \in \mathbb{R}^{80},
\]
as a normalized vector encoding global geometric information. One
possible construction is
\[
\tau_i(x)
    = \frac{1}{Z}
      \left(
        \sum_{j}
        \left| \nabla g_{ij}(x) \right|
      \right),
\]
where $g_{ij}(x)$ are the components of the metric tensor and $Z$ is a
normalization constant ensuring
\[
\|\tau(x)\|_\infty \leq 1.
\]

This signature reflects how the metric varies across the manifold and
thus captures global geometric behavior.

\subsection{Topology Modulation}

Given a holonomy-transported vector $h \in \mathbb{R}^{80}$, the
topology-modulated vector is defined as
\[
e = h \odot (1 + \beta \tau(x)),
\]
where $\beta$ is a scaling constant (typically $\beta = 0.05$ in the
Redux runtime).

This modulation incorporates global geometric information into the
representation.

\subsection{Stability Properties}

Because $\tau(x)$ depends smoothly on the metric tensor, we have
\[
\| \tau(x + \delta x) - \tau(x) \|
    \leq C \| \delta G \|,
\]
for some constant $C$ depending on the global geometric bounds of $M$.

Thus topology modulation preserves stability across the embedding
pipeline.

\subsection{Altitude Interpretation}

The topology operator completes the altitude A8 transformation by
incorporating global geometric structure. Whereas holonomy transport
captures accumulated curvature along canonical curves, topology
modulation captures how the manifold’s geometry varies globally.

\subsection{Output}

The output of the topology operator is the fully modulated embedding
vector
\[
e \in \mathbb{R}^{80},
\]
which serves as the input to the invariant computation layer in the
next stage of the Redux pipeline.
